\subsection{Sub-Phase 1.1: Geometric Representation via B-Splines and NURBS}
\label{subsec:nurbs_mapping}

The initial sub-phase was dedicated to building a high-performance geometric evaluator without external numerical dependencies, mirroring standard CAD kernels. This was achieved by implementing the Cox-de Boor recursive formulation to evaluate B-spline basis functions over an open uniform knot vector $\Xi = \{\xi_0, \xi_1, \dots, \xi_{m}\}$ where $m = n + p + 1$.

The algorithmic evaluation of a $p$-th degree basis function $N_{i,p}(\xi)$ relies on the classical recurrence relation:
\begin{equation}
N_{i,0}(\xi) = \begin{cases} 1 & \text{if } \xi_i \le \xi < \xi_{i+1} \\ 0 & \text{otherwise} \end{cases}
\end{equation}
\begin{equation}
N_{i,p}(\xi) = \frac{\xi - \xi_i}{\xi_{i+p} - \xi_i} N_{i,p-1}(\xi) + \frac{\xi_{i+p+1} - \xi}{\xi_{i+p+1} - \xi_{i+1}} N_{i+1,p-1}(\xi)
\end{equation}

For structural mechanics applications where variational strain energies involve higher-order spatial derivatives, the gradients of the basis functions are evaluated using the analytical derivative recurrence:
\begin{equation}
\frac{d}{d\xi} N_{i,p}(\xi) = p \left( \frac{N_{i,p-1}(\xi)}{\xi_{i+p} - \xi_i} - \frac{N_{i+1,p-1}(\xi)}{\xi_{i+p+1} - \xi_{i+1}} \right)
\end{equation}

To represent conic or non-isomorphic domains exactly, the basis was extended to non-uniform rational formulations ($\acrshort{nurbs}$) by injecting projective control weights $w_i$, defining the rational basis functions $R_{i,p}(\xi)$ as:
\begin{equation}
R_{i,p}(\xi) = \frac{N_{i,p}(\xi) w_i}{\sum_{j=0}^{n-1} N_{j,p}(\xi) w_j}
\end{equation}
The corresponding derivatives of the rational basis are computed using the quotient rule:
\begin{equation}
\frac{d}{d\xi} R_{i,p}(\xi) = w_i \frac{\frac{d N_{i,p}(\xi)}{d\xi} W(\xi) - N_{i,p}(\xi) \frac{d W(\xi)}{d\xi}}{[W(\xi)]^2}
\end{equation}
where $W(\xi) = \sum_{j=0}^{n-1} N_{j,p}(\xi) w_j$ is the denominator function representing the projective sum.

Unlike standard Lagrangian finite elements, NURBS basis functions do not possess the Kronecker delta property at internal nodes, meaning $R_{i,p}(\xi_j) \neq \delta_{ij}$. Consequently, a control point $\vect{P}_i$ does not represent the physical state at a specific coordinate unless it lies on a repeated boundary knot of multiplicity $p+1$ (clamped boundary), where the curve is forced to interpolate the control point directly. This requires advanced projection formulations for general boundary constraints. Additionally, NURBS shape functions offer $C^{p-1}$ continuity across knot spans, which significantly reduces numerical shear-locking and volume-locking compared to traditional $C^0$-continuous FEM elements.


\begin{figure}[H]
    \centering
    \includegraphics[width=0.75\textwidth]{nurbs_basis_curves.png}
    \caption{Visual representation of evaluated NURBS basis functions and their corresponding control weights extracted from the initial client-side prototype interface.}
    \label{fig:nurbs_curves_eval}
\end{figure}

The physical space coordinates of the parametric curve $C(\xi)$ are evaluated continuously by forming the linear combination of control points $\vect{P}_i = [x_i, y_i]^T$:
\begin{equation}
\vect{C}(\xi) = \sum_{i=0}^{n-1} R_{i,p}(\xi) \vect{P}_i
\end{equation}

\begin{devnote}
The early implementation inside \texttt{iga-sim-1.2/src/engine/iga.js} avoids full memory allocations by performing a direct, zero-dependency implementation of the Cox-de Boor recursive loop. To enforce standard physical conditions, the knot vectors are initialized as open uniform arrays, clamping the curve boundaries directly to the primary control points.
\end{devnote}